\chapter{Market structure and the limit order book}\label{ch:lob}
Chapters 1 and 2 treated prices as a given series. This chapter opens
the mechanism that prints them: the limit order book (LOB) --- who is
in it, how it matches, why a big order pays more than the quoted
price, and what data the venue emits so you can replay all of it.
Everything the harness in ch8 consumes was produced by the rules on
these pages.

\section{Who trades and why}\label{sec:whotrades}
\index{liquidity}\index{informed trader}\index{market maker}
Three kinds of flow share every book, and their incentives explain
the mechanics that follow:

\begin{itemize}
\item \textbf{Liquidity (uninformed) traders} want the position, not
the price --- they pay the spread to get filled now.
\item \textbf{Informed traders} know something about where the price
is going. They also cross the spread --- but only when their
information covers it.
\item \textbf{Market makers} sit on both sides, earning the spread
for standing ready. Their adversary is the second group: every fill
from someone who knows the near future is a loss waiting to be
marked. \S\ref{sec:makers} quantifies that loss.
\end{itemize}

The book is a queue discipline, not a negotiation. You either join a
queue (limit order, waits) or pay to jump it (market order, fills
now at worse prices). Everything else --- fees, rebates, order
types, even most ``signals'' --- is commentary on that choice.

\section{Orders and the LOB}\label{sec:orders}
\index{limit order book}\index{mid}\index{microprice}
The \emph{bid} is the best price someone will pay; the \emph{ask}
the best price someone will sell at. The \emph{spread} is their gap;
the \emph{mid} is its midpoint. \emph{Depth} is the quantity resting
at each level (Figure~\ref{fig:ladder}).

The mid is a convention, not a fact --- it treats both sides as
equally informative, and a book stacked on one side usually isn't.
The \emph{microprice} weights the quote by the size on the opposite
side:
\[
p_{\text{micro}}=\frac{a\,q_b + b\,q_a}{q_b+q_a},
\]
where $a,b$ are best ask/bid and $q_a,q_b$ their sizes.

\begin{example}[Mid vs microprice]\label{ex:microprice}
\texttt{code/microprice.py} on three book shapes. Real output:
\begin{lstlisting}[numbers=none]
balanced    book (4@100.0 | 2@100.5): mid=100.2500 microprice=100.3333
ask heavy   book (4@100.0 | 40@100.5): mid=100.2500 microprice=100.0455
bid heavy   book (40@100.0 | 2@100.5): mid=100.2500 microprice=100.4762
\end{lstlisting}
Same mid, three different honest mids: a wall on the ask side drags
the fair value down, a wall on the bid side pushes it up.
Microprice is the first ``information extraction'' computation of
the book --- imbalance before trade.
\end{example}

\begin{figure}[ht]\centering
\begin{tikzpicture}[x=1cm,y=1cm,
  lvl/.style={font=\small,anchor=east},
  bid/.style={fill=accentsoft,draw=accent,minimum width=2.6cm,minimum height=0.55cm},
  ask/.style={fill=amberbg,draw=amber,minimum width=2.6cm,minimum height=0.55cm}]
\node[font=\small\bfseries,text=slate] at (-1.2,3.3) {price};
\node[font=\small\bfseries,text=slate] at (1.4,3.3) {size};
\node[ask] (a3) at (1.4,2.6) {6}; \node[lvl] at (-1.2,2.6) {101.5};
\node[ask] (a2) at (1.4,1.9) {3}; \node[lvl] at (-1.2,1.9) {101.0};
\node[ask] (a1) at (1.4,1.2) {2}; \node[lvl] at (-1.2,1.2) {100.5};
\node[font=\small,text=slate,align=center] at (0.6,0.45)
  {spread $0.5$\\mid $100.25$};
\node[bid] (b1) at (1.4,-0.3) {4}; \node[lvl] at (-1.2,-0.3) {100.0};
\node[bid] (b2) at (1.4,-1.0) {3}; \node[lvl] at (-1.2,-1.0) {99.5};
\node[bid] (b3) at (1.4,-1.7) {5}; \node[lvl] at (-1.2,-1.7) {99.0};
\node[font=\small,text=accent,anchor=west] at (3.2,-0.65) {bids: high first};
\node[font=\small,text=amber,anchor=west] at (3.2,2.25) {asks: low first};
\draw[<->,slate] (-0.6,1.2) -- (-0.6,-0.3);
\end{tikzpicture}

\caption{The order-book ladder. Asks sorted low-to-high above the
spread, bids high-to-low below it. Depth is the queue behind each
price.}
\label{fig:ladder}
\end{figure}

\begin{example}[Price--time priority, by hand]\label{ex:match}
Into the Figure~\ref{fig:ladder} book goes a limit \emph{sell} of 5
at 99.5 (\texttt{code/lob.py}). It is marketable: bids exist at or
above 99.5. Real output:
\begin{lstlisting}[numbers=none]
top: (100.0, 100.5) mid: 100.25
bids: [(100.0, 4), (99.5, 3), (99.0, 5)]
asks: [(100.5, 2), (101.0, 3), (101.5, 6)]
after sell 5@99.5: [(100.0, 4, 's'), (99.5, 1, 's')]
top: (99.5, 100.5)
\end{lstlisting}
The seller did not choose who filled --- the book did: best price
first (100.0), then time order within that level, and the leftover
unit rested as the new best ask. Two lessons: a limit sell below
the top bid \emph{trades immediately} (a limit order is ``pay up to
here'', not ``exactly here''), and the new best \emph{ask} is 99.5
because the seller's remainder now sits there.
\end{example}

\begin{checkpoint}[Which side trades?]
A limit buy at 101.0 arrives against the
Figure~\ref{fig:ladder} book. How many units fill, at what
prices, and what rests? (5 fill --- 2@100.5 then 3@101.0; nothing
rests; new ask top is 101.5.)
\end{checkpoint}

\section{Matching: walking the book is the slippage}\label{sec:match}
\index{slippage}\index{price--time priority}
A market order does not get ``the price'' --- it gets the queue.
Each level it eats is paid that level's price, so the average fill
is a volume-weighted walk up the ladder, and the gap to mid is
\emph{slippage}.

\begin{example}[A market order through the ladder]\label{ex:walk}
Market buy 9 against asks $(100.5\!\times\!2,\;101.0\!\times\!3,\;
101.5\!\times\!6,\;102.5\!\times\!10)$ ---
\texttt{code/walk\_book.py}. Real output:
\begin{lstlisting}[numbers=none]
fills: [(100.5, 2, 'b'), (101.0, 3, 'b'), (101.5, 4, 'b')]
filled 9 at avg 101.1111  vs mid 100.25
slippage = +0.8611/unit = +0.859%
posted spread was 0.5; paid avg above best ask by +0.6111
\end{lstlisting}
The poster's ``cost'' is not the 0.5 spread --- it is the average
$+0.86$ above mid, convex in size (Figure~\ref{fig:slip}): each
level eaten costs more than the last, and a 9-unit order never saw
the quoted 2-unit price except for its first 2 units.
\end{example}

\begin{figure}[ht]\centering
\begin{tikzpicture}
\begin{axis}[width=0.7\linewidth,height=5.2cm,
  xlabel={market-buy size (units)},ylabel={slippage vs mid (\%)},
  grid=major,grid style={gray!25},xmin=0,xmax=48,ymin=0,ymax=3.5]
\addplot[accent,thick,mark=*] coordinates {(2,0.249) (3,0.416) (5,0.549) (6,0.665) (9,0.859) (12,1.039) (21,1.556) (36,2.466) (46,3.177)};
\node[font=\small,text=slate,anchor=west] at (axis cs:6,2.9)
  {convex: each level eaten costs more than the last};
\end{axis}
\end{tikzpicture}

\caption{Depth$\to$slippage on the Figure~\ref{fig:ladder} ladder
extended with deeper levels: average-fill distance from mid grows
convexly in order size.}
\label{fig:slip}
\end{figure}

\begin{figure}[ht]\centering
\begin{tikzpicture}[font=\small,
  ask/.style={fill=amberbg,draw=amber,minimum width=1.7cm,minimum height=0.6cm},
  eat/.style={fill=amber!45!white,draw=amber,minimum width=1.7cm,minimum height=0.6cm},
  lbl/.style={font=\footnotesize,text=slate}]
\foreach \x/\t in {0/{book before},3.4/{fill 2@100.5},6.8/{fill 3@101.0},10.2/{fill 4@101.5}}
  \node[font=\footnotesize\bfseries,text=slate] at (\x+0.85,3.15) {\t};
\node[ask] at (0.85,2.4) {6@101.5};
\node[ask] at (0.85,1.7) {3@101.0};
\node[ask] at (0.85,1.0) {2@100.5};
\node[lbl] at (0.85,0.3) {market buy 9};
\node[eat] at (4.25,1.0) {2@100.5};
\draw[->,accent,thick] (2.0,0.3) -- (3.3,0.8);
\node[ask] at (4.25,2.4) {6@101.5};
\node[ask] at (4.25,1.7) {3@101.0};
\node[eat] at (7.65,1.7) {3@101.0};
\node[ask] at (7.65,2.4) {6@101.5};
\draw[->,accent,thick] (5.3,0.9) -- (6.6,1.4);
\node[eat] at (11.05,2.4) {4@101.5};
\node[ask] at (11.05,3.0) {2@101.5 left};
\draw[->,accent,thick] (8.7,1.6) -- (10.0,2.1);
\node[lbl,align=left] at (11.0,0.5) {avg fill $101.11$\\slippage $+0.86$ vs mid};
\end{tikzpicture}

\caption{The same walk, drawn: the market buy consumes levels
cheapest-first and the unfilled remainder of the top level keeps
resting.}
\label{fig:walk}
\end{figure}

\begin{figure}[ht]\centering
\begin{tikzpicture}[font=\small,
  st/.style={draw=slate,rounded corners,fill=graybg,minimum height=0.7cm,align=center},
  ev/.style={->,accent,thick},
  lbl/.style={font=\footnotesize,text=slate,fill=white,inner sep=1.5pt}]
\node[st] (open) at (0,0) {open order\\(resting)};
\node[st] (book) at (4.6,0) {in the book\\(queued by time)};
\node[st] (fill) at (9.2,0) {filled};
\node[st] (cxl)  at (4.6,-1.8) {cancelled};
\draw[ev] (open) -- node[lbl,above]{not marketable} (book);
\draw[ev] (book) -- node[lbl,above]{counterparty order arrives} (fill);
\draw[ev] (book) -- node[lbl,right]{owner cancels / IOC timeout} (cxl);
\draw[ev] (open) edge[bend left=25] node[lbl,above]{marketable: immediate match} (fill);
\draw[ev] (open) edge[bend right=15] node[lbl,below]{cancel / expire} (cxl);
\end{tikzpicture}

\caption{One resting order's life. Marketable orders skip the book
entirely; everything else queues by price then time.}
\label{fig:state}
\end{figure}

\section{Makers, takers, fees and adverse selection}\label{sec:makers}
\index{maker/taker}\index{adverse selection}
Venues charge the two sides differently: \emph{takers} (marketable
orders, consume liquidity) pay more; \emph{makers} (resting orders,
provide it) pay less --- sometimes nothing, sometimes a rebate.
Figure~\ref{fig:fees} shows one venue's base tiers as an example;
treat every such number as a dated claim, because tiers change and
yours will not be the base.

\begin{figure}[ht]\centering
\begin{tikzpicture}
\begin{axis}[width=0.72\linewidth,height=4.6cm,ybar,bar width=0.55cm,
  ylabel={fee (bp of notional)},symbolic x coords={maker,taker},
  xtick=data,ymin=0,ymax=6,grid=major,grid style={gray!25},
  nodes near coords,nodes near coords style={font=\small}]
\addplot[fill=accent,draw=none] coordinates {(maker,1.5) (taker,4.5)};
\end{axis}
\node[font=\footnotesize,text=slate,align=left] at (5.6,-0.4)
  {Hyperliquid perp base tiers, docs as of 2026-10-10 [unverified --- re-check]};
\end{tikzpicture}

\caption{Maker/taker fee asymmetry at one venue's base tier
(Hyperliquid perp documentation, accessed 2026-10-10
\textbf{[unverified --- re-check current tiers]}; volume tiers reduce
both). Order-type names and fee schedules differ per venue; on
Polymarket the CLOB displays prices as probabilities with its own
midpoint display rule (docs accessed 2026-10-10 \textbf{[unverified]}).}
\label{fig:fees}
\end{figure}

Why makers get paid at all: their real cost isn't the fee, it is
\emph{adverse selection} --- the fills that come from people who
know the next move.

\begin{example}[Adverse selection, measured]\label{ex:mmgame}
\texttt{code/mm\_game.py}: a maker quotes $\pm0.10$ around fair
value for 3000 steps. Uninformed flow coin-flips; informed flow
trades only when the next move covers the spread. 24 seeds per
flow mix. Real output:
\begin{lstlisting}[numbers=none]
informed share 0.0: mean MM PnL   +298.51
informed share 0.3: mean MM PnL   +191.02 | per informed trade -0.018
informed share 0.6: mean MM PnL    +97.74 | per informed trade -0.019
informed share 0.9: mean MM PnL    +39.17 | per informed trade -0.019
\end{lstlisting}
Same quotes, same spread --- only the flow changed. Every informed
trade is a guaranteed loss of (move $-$ spread); the maker's entire
job is widening or withdrawing quotes when informed share rises.
This is why real quoting engines reprice on flow toxicity, not on
the mid alone.
\end{example}

\begin{checkpoint}[Where did the edge go?]
In Example~\ref{ex:mmgame} the maker still nets positive at 90\%
informed. Why, and what two quote changes flip it negative?
(Uninformed flow still pays $+0.10$ per trade; widening the spread
or pulling quotes during informed bursts starves the adverse
fills.)
\end{checkpoint}

\section{What the tape looks like}\label{sec:tape}
\index{trades}\index{book deltas}
A venue emits three streams: \emph{trades} (prints: price, size,
taker side, time), \emph{book deltas} (incremental
add/change/delete per level), and derived \emph{candles} (OHLC
aggregates). Ch1's bars are the third stream; the first two are the
raw material it is built from --- and you can rebuild it yourself to
check.

\begin{example}[Rebuilding candles from real prints]\label{ex:candles}
Fetch $\sim$1000 real BTC trade prints from Kraken's public Trades
endpoint, resample into 60-second candles, and compare with the
venue's own 1-minute OHLC for the same window
(\texttt{code/trades\_to\_candles.py}; endpoints verified
2026-10-11, cached in \texttt{data/btc\_trades.csv}). Real output,
aligned interior minutes:
\begin{lstlisting}[numbers=none]
resampled: t=...020 n=132 o=83812.8 h=83846.6 l=83812.7 c=83843.4
venue:     t=...020       o=83812.8 h=83846.6 l=83812.7 c=83843.4
resampled: t=...080 n=170 o=83843.4 h=83873.9 l=83843.3 c=83860.6
venue:     t=...080       o=83843.4 h=83873.9 l=83843.3 c=83860.6
resampled: t=...140 n=106 o=83857.9 h=83880.0 l=83853.6 c=83876.3
venue:     t=...140       o=83857.9 h=83880.0 l=83853.6 c=83876.3
\end{lstlisting}
Every complete minute matches to the tick. The edge minutes differ
--- the print window is partial there --- which is itself the lesson:
candles are lossy views of a print stream; reconstruct before you
trust. The buy/sell mix in the same fetch (748 taker buys vs 258
sells) also shows how one-sided real flow is ---
Figure~\ref{fig:bounce}.
\end{example}

\begin{figure}[ht]\centering
\begin{tikzpicture}
\begin{axis}[width=0.86\linewidth,height=5.4cm,
  xlabel={trade \#},ylabel={price (demeaned)},
  grid=major,grid style={gray!25},xmin=0,xmax=160,
  legend style={font=\small}]
\addplot[only marks,mark=*,mark size=1.2pt,color=accent] coordinates {(0,9.74) (1,9.74) (2,9.74) (7,7.94) (10,7.64) (11,7.64) (32,-0.46) (33,-0.46) (39,-4.66) (40,-4.66) (49,-4.66) (50,-4.66) (51,-4.66) (52,-4.66) (53,-4.66) (60,-15.26) (62,-15.26) (63,-15.26) (64,-13.26) (65,-10.96) (66,-10.86) (67,-10.86) (68,-9.06) (70,-9.06) (71,-9.06) (72,-9.06) (73,-9.06) (74,-9.06) (76,-9.06) (77,-9.06) (78,-9.06) (79,-9.06) (80,-9.06) (81,-9.06) (82,-4.86) (83,-2.36) (84,-0.66) (85,7.64) (86,7.74) (87,9.24) (88,9.24) (90,9.24) (92,9.24) (93,9.24) (94,9.24) (95,9.24) (96,9.24) (97,9.24) (98,9.24) (99,9.24) (100,9.24) (101,9.24) (102,9.24) (103,9.24) (114,5.04) (115,5.04) (124,-1.16) (125,-1.16) (126,-1.16) (127,-1.16) (128,-1.16) (129,-1.16) (130,-1.16) (131,-1.16) (132,-1.16) (133,-1.16) (134,-1.16) (135,-1.16) (136,-1.16) (137,-0.66) (138,0.94) (139,0.94) (140,0.94) (141,0.94) (142,3.54) (143,5.04) (149,3.54) (154,3.54) (156,3.54) (157,3.54) (158,3.54) (159,3.54)};
\addlegendentry{taker buy (lifts ask)}
\addplot[only marks,mark=square*,mark size=1.2pt,color=amber] coordinates {(3,9.64) (4,9.64) (5,9.64) (6,9.64) (8,7.84) (9,7.84) (12,7.54) (13,7.54) (14,7.54) (15,7.54) (16,3.64) (17,-0.56) (18,-4.56) (19,-4.56) (20,-4.76) (21,-7.66) (22,-8.96) (23,-10.86) (24,-10.86) (25,-11.86) (26,-11.96) (27,-12.06) (28,7.54) (29,7.54) (30,7.54) (31,3.64) (34,-0.56) (35,-0.56) (36,-0.56) (37,-0.56) (38,-0.56) (41,-4.76) (42,-4.76) (43,-8.96) (44,-13.16) (45,-13.96) (46,-13.96) (47,-16.66) (48,-4.76) (54,-4.76) (55,-4.76) (56,-4.76) (57,-8.96) (58,-13.16) (59,-15.46) (61,-15.36) (69,-9.16) (75,-9.16) (89,9.14) (91,9.14) (104,9.14) (105,9.14) (106,9.14) (107,9.14) (108,9.14) (109,9.14) (110,7.84) (111,3.64) (112,2.84) (113,2.74) (116,4.94) (117,4.94) (118,4.94) (119,4.94) (120,4.94) (121,3.64) (122,-0.56) (123,-1.26) (144,4.94) (145,4.94) (146,4.94) (147,3.64) (148,-0.36) (150,3.44) (151,3.44) (152,3.44) (153,3.44) (155,3.44)};
\addlegendentry{taker sell (hits bid)}
\addplot[ink,densely dotted,thick] coordinates {(3,9.69) (4,9.69) (5,9.69) (6,9.69) (7,8.79) (8,7.89) (9,7.89) (10,7.74) (11,7.74) (12,7.59) (13,7.59) (14,7.59) (15,7.59) (16,5.64) (17,3.54) (18,1.54) (19,1.54) (20,1.44) (21,-0.01) (22,-0.66) (23,-1.61) (24,-1.61) (25,-2.11) (26,-2.16) (27,-2.21) (28,7.59) (29,7.59) (30,7.59) (31,5.64) (32,1.59) (33,1.59) (34,-0.51) (35,-0.51) (36,-0.51) (37,-0.51) (38,-0.51) (39,-2.61) (40,-2.61) (41,-4.71) (42,-4.71) (43,-6.81) (44,-8.91) (45,-9.31) (46,-9.31) (47,-10.66) (48,-4.71) (49,-4.71) (50,-4.71) (51,-4.71) (52,-4.71) (53,-4.71) (54,-4.71) (55,-4.71) (56,-4.71) (57,-6.81) (58,-8.91) (59,-10.06) (60,-15.36) (61,-15.31) (62,-15.31) (63,-15.31) (64,-14.31) (65,-13.16) (66,-13.11) (67,-13.11) (68,-12.21) (69,-9.11) (70,-9.11) (71,-9.11) (72,-9.11) (73,-9.11) (74,-9.11) (75,-9.11) (76,-9.11) (77,-9.11) (78,-9.11) (79,-9.11) (80,-9.11) (81,-9.11) (82,-7.01) (83,-5.76) (84,-4.91) (85,-0.76) (86,-0.71) (87,0.04) (88,0.04) (89,9.19) (90,9.19) (91,9.19) (92,9.19) (93,9.19) (94,9.19) (95,9.19) (96,9.19) (97,9.19) (98,9.19) (99,9.19) (100,9.19) (101,9.19) (102,9.19) (103,9.19) (104,9.19) (105,9.19) (106,9.19) (107,9.19) (108,9.19) (109,9.19) (110,8.54) (111,6.44) (112,6.04) (113,5.99) (114,3.89) (115,3.89) (116,4.99) (117,4.99) (118,4.99) (119,4.99) (120,4.99) (121,4.34) (122,2.24) (123,1.89) (124,-1.21) (125,-1.21) (126,-1.21) (127,-1.21) (128,-1.21) (129,-1.21) (130,-1.21) (131,-1.21) (132,-1.21) (133,-1.21) (134,-1.21) (135,-1.21) (136,-1.21) (137,-0.96) (138,-0.16) (139,-0.16) (140,-0.16) (141,-0.16) (142,1.14) (143,1.89) (144,4.99) (145,4.99) (146,4.99) (147,4.34) (148,2.34) (149,1.59) (150,3.49) (151,3.49) (152,3.49) (153,3.49) (154,3.49) (155,3.49) (156,3.49) (157,3.49) (158,3.49) (159,3.49)};
\addlegendentry{approx mid}
\end{axis}
\end{tikzpicture}

\caption{Real Kraken prints, first 160 of the cached window: taker
buys print at the ask, sells at the bid, and the inferred mid
bounces between. Even in a trending window the two-side split is
visible.}
\label{fig:bounce}
\end{figure}

\begin{example}[Measuring signed impact on real prints]\label{ex:signimpact}
\textbf{Strategy.} With no quotes pinned, proxy the half-spread by
how far the price moves after a buy vs after a sell. On the pinned
burst, split consecutive price changes by the sign of the
\emph{previous} print.

\textbf{Solution.} \texttt{signed\_impact.py} prints:
\begin{verbatim}
prints 1006: buys 748 sells 258
zero-move prints: 54% of consecutive moves are 0
avg next-print dp after a buy  = +0.3760
avg next-print dp after a sell = -0.5101
signed impact (spread/2 proxy) = 0.4431 per print
avg |dp| = 0.655 = 6.55 ticks
\end{verbatim}
Buys push the next print up on average, sells push it down --- that
is signed price impact, mechanically identical to the adverse
selection the MM felt in Example~\ref{ex:mmgame}. The proxy
half-spread $\approx 0.44$ per print ($\approx 4.4$ ticks) is far
tighter than a naive ``last price'' read because 54\% of moves are
zero --- most prints re-mark at the same level.

\textbf{Significance.} You just estimated a transaction-cost number
from raw trades alone: when quotes are missing, signed impact is
the cheapest workable spread proxy. It is biased downward (it sees
only \emph{realised} moves, not the queue), which is why real TCA
keeps the book snapshots.
\end{example}

\section{A 60-line exchange simulator}\label{sec:sim}
\index{simulator}
\texttt{code/lob.py} is a working exchange in one file: price-level
queues, price--time priority, immediate matching for marketable
orders, resting for the rest --- everything in this chapter's WEs ran
on it. Read it as the harness preview: ch8's backtest engine is this
same event loop (arrive $\to$ match $\to$ mark) with fill realism,
latency and the no-lookahead rule added. If you can explain every
line of \texttt{lob.py}, you already understand what a backtester
does all day.

\section*{Summary sheet}
\addcontentsline{toc}{section}{Summary sheet}
\begin{table}[ht]\centering
\begin{tabularx}{\linewidth}{@{}lYY@{}}
\toprule
\tablehead\thead{Object} & \thead{Definition} & \thead{When to use}\\
\midrule
Spread / mid & $a-b$, $(a+b)/2$ & quoting context, not a fair price\\
Microprice & $(a\,q_b+b\,q_a)/(q_a+q_b)$ & imbalance-informed mid\\
Price--time & best price, then FIFO & who fills first\\
Slippage & avg fill $-$ mid & real cost of market orders\\
Maker/taker & rest vs cross the spread & fee + adverse-selection
trade-off\\
Prints / deltas / candles & the three venue streams & raw data
before aggregates\\
\bottomrule
\end{tabularx}
\caption{Chapter 3 toolkit.}
\end{table}

\begin{example}[Order-flow imbalance predicts size, not side]\label{ex:ofi}
\textbf{Strategy.} Without a book feed, reconstruct the closest
thing to queue pressure: order-flow imbalance (signed volume in a
rolling 100-print window) and test whether it predicts the next
print's move in sign or in magnitude.

\textbf{Solution.} \texttt{book\_regen.py} prints:
\begin{verbatim}
OFI vs next-print move over 1006 prints
corr(OFI, dp_next) = -0.040
mean |dp| when |OFI/vol|>0.5: 1.1680
mean |dp| when |OFI/vol|<0.1: 0.9187
\end{verbatim}
Direction: nothing ($r=-0.04$). Magnitude: imbalanced flow
precedes $27\%$ larger next-print moves --- the imbalance is
energy, not a compass.

\textbf{Significance.} This is the honest version of ``order flow
predicts price'': at print granularity the signal is volatility,
not drift. A model that trades the \emph{sign} of OFI buys noise;
one that sizes around it is reading the tape correctly. It also
validates the chapter's claim that microstructure edges need
book-level state --- prints alone tell you only how hard it will
move, not where.
\end{example}

\begin{example}[Fast fills are expensive fills]\label{ex:queue}
\textbf{Strategy.} Your resting bid sits behind $60$ queued units.
The queue drains two ways --- market sells eat the volume ahead of
you, and your queue-mates cancel. Simulate time-to-fill under calm,
busy and panic arrival rates, then mark out the mid after each
fill: does a fast fill arrive with a bill?

\textbf{Solution.} \texttt{queue\_sim.py} prints:
\begin{verbatim}
 calm: fill prob 1.00  median wait  27.1s  exp markout   +0.0bp
 busy: fill prob 1.00  median wait  11.5s  exp markout  -11.6bp
panic: fill prob 1.00  median wait   3.5s  exp markout  -35.4bp
faster fills are not free: they arrive exactly when informed flow is trading through you
\end{verbatim}
Wait time collapses $27\to3.5$s as flow intensifies, but the
expected markout goes from $0$ to $-35$bp: the fills that come
fastest are the ones informed sellers chose to give you.

\textbf{Significance.} Queue mechanics convert two rates ---
trade arrival and cancellation --- into a personal fill clock, and
the clock runs fastest precisely when you least want the fill.
This is the maker's dilemma in one table: your fill probability
and your adverse-selection bill are the same variable. A quote's
value is not set by the level's depth but by \emph{your rank} in it
and by \emph{who} is consuming the queue.
\end{example}

\section{Common mistakes}\index{common mistakes}
\begin{itemize}
\item \textbf{Costing a market order as ``the spread''.} You pay
spread once on top-of-book; size past the first level pays the
ladder (Example~\ref{ex:walk}).
\item \textbf{Treating mid as fair value.} Mid ignores imbalance;
use microprice or better when the book is lopsided
(Example~\ref{ex:microprice}).
\item \textbf{Assuming a limit order fills at its limit.} It fills
at the best available price --- a limit sell below the top bid
trades at the bid (Example~\ref{ex:match}).
\item \textbf{Forgetting your order joins a queue.} Behind 1000
units at your price, the level can trade hundreds before you.
\item \textbf{Trusting venue candles blindly.} They are aggregates;
the print stream is the record --- rebuild and compare
(Example~\ref{ex:candles}).
\item \textbf{Quoting without a flow model.} A static spread earns
$+0.10$ from noise traders and loses it back to informed flow
(Example~\ref{ex:mmgame}).
\item \textbf{Trusting ``last price'' as value.} 54\% of
consecutive prints in the pinned burst mark no move
(Example~\ref{ex:signimpact}) --- the last traded price is stale the
moment it prints. Mid or microprice is the honest mark.
\item \textbf{Forgetting impact has a sign.} Buys push up, sells
push down, \emph{by construction}. A backtest that fills at last
price assumes you moved the market $0.44$ ticks per trade in your
favour --- it gives you the rebate you never earned.
\end{itemize}

\section*{Exercises}\addcontentsline{toc}{section}{Exercises}
\begin{exercise}[easy]
Book: bids $(99.8\!\times\!3, 99.7\!\times\!2)$, asks
$(100.1\!\times\!2, 100.3\!\times\!4)$. Compute spread, mid,
microprice.
\end{exercise}
\begin{solution}
Spread $0.3$, mid $99.95$, microprice
$(100.1\!\cdot\!3+99.8\!\cdot\!2)/5=99.98$.
\end{solution}
\begin{exercise}[easy]
A market sell of 4 hits bids $(100.0\!\times\!2,99.9\!\times\!5)$.
Average fill and slippage vs mid $100.15$?
\end{exercise}
\begin{solution}
$(2\!\cdot\!100.0+2\!\cdot\!99.9)/4=99.95$; slippage
$99.95-100.15=-0.20$ per unit.
\end{solution}
\begin{exercise}[easy]
State price--time priority and what a limit order priced past the
best opposite level does.
\end{exercise}
\begin{solution}
Best price first, FIFO within a price; a marketable limit order
fills immediately at better prices up to its limit.
\end{solution}
\begin{exercise}[easy]
Define maker and taker, and say which typically pays the higher
fee and why the venue prices it that way.
\end{exercise}
\begin{solution}
Taker crosses the spread (consumes liquidity) and pays more; maker
rests (provides liquidity) and is subsidised because depth is what
the venue sells.
\end{solution}
\begin{exercise}[easy]
Name the three data streams a venue emits and which one ch1's
candles come from.
\end{exercise}
\begin{solution}
Trades (prints), book deltas, and derived candles; candles are the
derived aggregate of the print stream.
\end{solution}
\begin{exercise}[medium]
Effective vs posted spread: a taker buys 3, paying 2@100.5 and
1@101.0 against mid 100.25. Compute the effective spread (twice the
signed distance from mid per unit) and compare with posted 0.5.
\end{exercise}
\begin{solution}
Avg $100.667$; effective $=2\times(100.667-100.25)=0.83$ vs posted
$0.5$ --- walking depth inflates the realized cost.
\end{solution}
\begin{exercise}[medium]
Why does the mid move after fills even when no new quotes arrive?
Trace one in Example~\ref{ex:match}.
\end{exercise}
\begin{solution}
A fill can delete or demote the best level: the sell of 5 emptied
100.0's queue and left a 99.5 ask, so the top became
$(99.5,100.5)$ and mid fell from 100.25 to 100.0 --- trades
\emph{are} quote updates.
\end{solution}
\begin{exercise}[medium]
Queue position: you place a limit buy at the current best bid,
behind 400 units. A 250-unit sell market order arrives, then a
200-unit one. Do you fill? How much?
\end{exercise}
\begin{solution}
The queue ahead absorbs 400 of the 450 sold; you fill 50 --- FIFO
within the level, no pro-rata in a price--time book.
\end{solution}
\begin{exercise}[medium]
Extend \texttt{lob.py} with a \texttt{cancel(side, price)} that
removes the front order at a level. What does the deque need to
support to cancel a \emph{specific} order id instead?
\end{exercise}
\begin{solution}
Front-cancel is \texttt{popleft}; cancelling an arbitrary id needs
either per-id index or tombstoning --- remove by marking, skip on
match.
\end{solution}
\begin{exercise}[medium]
In \texttt{trades\_to\_candles.py}, why do the edge minutes differ
between resampled and venue candles while interiors match exactly?
\end{exercise}
\begin{solution}
The fetched print window starts/ends mid-minute, so those buckets
are partial; interior minutes contain every print in range, and OHLC
of the same complete set is deterministic.
\end{solution}
\begin{exercise}[hard]
Adverse-selection toy: with spread $S$, per-step move sd $\sigma$,
and informed share $\eta$ trading only when $|move|>S$, write the
expected per-trade maker PnL and find the $\eta$ where it crosses
zero at $S=0.10$, $\sigma=0.05$.
\end{exercise}
\begin{solution}
Uninformed pays $+S$ per trade (50\% of orders); informed costs
$E[m-S\mid m>S]\approx$ a few $\sigma/100$. Break-even when informed
flow's expected drag equals the uninformed capture --- numerically
$\eta\gtrsim0.9$ in Example~\ref{ex:mmgame}'s run, sensitive to
$\sigma/S$.
\end{solution}
\begin{exercise}[hard]
Design the smallest change to \texttt{mm\_game.py} that lets the
maker \emph{adapt}: which observable (hint: it is in every trade
print) predicts adverse flow, and what quote response cuts the loss?
\end{exercise}
\begin{solution}
Taker-side imbalance/flow toxicity: track recent buy vs sell share
and skew quotes (widen the hit side, shade mid) or widen both when
one-sided flow persists --- the queue then prices the information,
not just the mid.
\end{solution}
\begin{exercise}[hard]
A venue offers a maker rebate of $-1$bp. Your quotes fill 60\% of
the time and adversely select 0.5bp on average. Is quoting
profitable? What else must you check before concluding?
\end{exercise}
\begin{solution}
Per fill you net $+1$bp rebate $-$ $0.5$bp adverse $=+0.5$bp gross
--- but only if fills are random; check fill\emph{rate asymmetry}
(you fill most when you're wrong), inventory cost, and fees on the
hedging leg before concluding.
\end{solution}
\begin{exercise}[hard]
Reconstruct \texttt{lob.py}'s matching as events: write the exact
sequence of book mutations for a marketable limit sell of 5@99.5
arriving at the Figure~\ref{fig:ladder} book (fill events and the
residual add), then run it and diff against your trace.
\end{exercise}
\begin{solution}
Fill(100.0,$\times$4$\to$0, delete level), fill(99.5,$\times$1 of 3
$\to$ 2 left), residual 0 --- the output shows two fills and a new
best ask 99.5... wait, verify by running: the trace must match
\texttt{lob.py}'s real output exactly.
\end{solution}
\begin{exercise}[hard]
Implement an effective-spread and a realized-spread function (fill
list, mid, and a mark-out mid at $t{+}60$s); explain why the second
splits adverse selection from liquidity cost.
\end{exercise}
\begin{solution}
Effective $=2|\bar p-m_0|$; realized $=2\,\mathrm{sign}\cdot
(\bar p-m_{t+60})$ --- the part that stays after the mark-out window
is real cost, the part that reverts was informed-flow noise.
\end{solution}

\begin{exercise}[easy]
In Example~\ref{ex:signimpact} the signed-impact proxy is
$0.44$/print. Is that the spread, the half-spread, or neither?
Justify in one line.
\end{exercise}
\begin{solution}
A half-spread proxy: it measures one-sided impact per print
(buy-vs-sell gap divided by two), i.e.\ what one side of the
round trip costs --- not the full bid-ask width.
\end{solution}
\begin{exercise}[medium]
In the burst, buys outnumber sells 748 to 258 yet price still finds
sellers. Propose two mechanisms a live book provides that the print
file hides.
\end{exercise}
\begin{solution}
(a) Resting limit sells absorb buy aggression without printing
``sell'' --- the tape shows the taker's side only; (b) market
makers re-quote continuously, so the ask refills between prints.
Print counts measure taker aggression, not liquidity.
\end{solution}
\begin{exercise}[hard]
Design an estimator of the \emph{effective} spread using only the
print file and a proxy mid built from a smoothed (e.g.\ 5-print
median) price. What biases does the smoothed mid introduce vs a
true quote mid?
\end{exercise}
\begin{solution}
Eff.\ spread $= 2\cdot E[\,s_t(p_t - \tilde m_t)\,]$ with
$\tilde m$ the smoothed mid. Bias: the smoothed mid lags fast moves
so buy/sell classifications get noisier --- it understates the
spread in trends, overstates in chop.
\end{solution}

\section*{Notes}\addcontentsline{toc}{section}{Notes}
\textbf{Venue rules are dated claims.} Fee tiers, order-type names
and display rules cited here (Hyperliquid, Polymarket) were read
from venue docs on 2026-10-10 and are \textbf{[unverified]} --- see
\texttt{docs/unverified.md}; re-check before relying.
\textbf{Data.} \texttt{code/data/btc\_trades.csv} is a pinned Kraken
Trades snapshot (fetched 2026-10-11); \texttt{btc\_daily.csv} covers
ch1--2. \textbf{Standing disclaimer.} Research methods and market
mechanics only; no investment, legal or tax advice.
\textbf{Where this leads.} Ch4 adds the funding instrument that
prices the dominant venue's perpetuals; ch10 prices the fills this
chapter walked.

\section*{Further reading}\addcontentsline{toc}{section}{Further reading}
\begin{itemize}
\item Polymarket developer docs --- a live CLOB with prices as
probabilities; the midpoint-display rule cited in \S\ref{sec:makers}:

\url{https://docs.polymarket.com/}

\item Hyperliquid docs --- the venue whose base fee tiers appear in
Figure~\ref{fig:fees} (accessed 2026-10-10, tiers
\textbf{[unverified]}):

\url{https://hyperliquid.gitbook.io/hyperliquid-docs/trading/fees}

\item Kraken market-data REST docs --- the Trades and OHLC endpoints
behind Example~\ref{ex:candles}:

\url{https://docs.kraken.com/api-reference/market-data/get-recent-trades}

\item M.~L\'opez de Prado, lecture library --- microstructure and
adverse-selection lectures beyond this chapter:

\url{https://www.quantresearch.org/Lectures.htm}

\item W.~McKinney, \emph{Python for Data Analysis} --- the resampling
idioms behind the candle-rebuild example:

\url{https://wesmckinney.com/book/}
\end{itemize}

\begin{rulebox}[Watch instead]
\begin{itemize}
\item \textbf{Initial Return --- ``Limit order book / Fundamentals
of trading / Episode 4''}, 18:56 --- inside quote, mid, price--time
priority and round-trip cost; maps to
\S\ref{sec:orders}--\S\ref{sec:match}. (Verified 2026-10-11.)
\url{https://www.youtube.com/watch?v=BJt6PGKeO9I}
\item \textbf{INFORMS --- ``Stochastic Market Microstructure Models
of Limit Order Books''}, 1:28:22 --- the stochastic treatment behind
queue dynamics and adverse selection; maps to
\S\ref{sec:whotrades} and \S\ref{sec:makers}. (Verified 2026-10-11.)
\url{https://www.youtube.com/watch?v=XoBjQqMmKoM}
\end{itemize}
\end{rulebox}
